\section{Enfoque competencial: niveles de desempeño y evidencias}
La evaluación por competencias es el eje de la propuesta: cada actividad genera una evidencia y cada evidencia se valora con un nivel de desempeño. La tabla define cuatro niveles por competencia específica; las rúbricas de la sección~\ref{sec:instrumentos} los concretan.

\begin{tablaSAlarga}{colspec={Q[l,wd=30pt]X[l]X[l]X[l]X[l]Q[l,wd=74pt]},
  cells={font=\scriptsize}, colsep=3pt,
  row{1}={bg=urjc, fg=white, font=\scriptsize\bfseries},
  column{1}={font=\scriptsize\bfseries, fg=urjc}, cell{1}{1}={fg=white}}
  CE & 1 · Inicial & 2 · En desarrollo & 3 · Adecuado & 4 · Excelente & Evidencias \\
  CE1 & Aplica fórmulas sin relacionarlas con el problema & Modeliza con ayuda; omite alguna solución & Modeliza de forma autónoma y obtiene todas las soluciones & Además compara estrategias y justifica la más eficiente & Libro, F3, examen \\
  CE2 & No comprueba sus resultados & Comprueba con GeoGebra sin explicar diferencias & Comprueba y explica las discrepancias & Además estima el efecto del error y elige la solución según el contexto & F1, libro \\
  CE3 & Reproduce resultados dados & Formula conjeturas con guía, sin justificarlas & Formula conjeturas y las demuestra con ayuda & Demuestra de forma autónoma y propone casos nuevos & F2, examen \\
  CE5 & Trata cada contenido por separado & Relaciona dos ideas con ayuda & Conecta vectores, ángulos y rectas & Usa las conexiones para resolver de varias formas & Mapa, examen \\
  CE6 & No relaciona matemáticas y entorno & Señala alguna conexión general & Aplica el modelo a la instalación real & Además valora impacto y límites del modelo & Libro, defensa \\
  CE7 & Representaciones incorrectas o incompletas & Representa con una sola forma & Selecciona la representación adecuada & Justifica la elección y comunica con claridad & F1, F4, libro \\
  CE8 & Lenguaje impreciso y desordenado & Comunica con errores de notación & Comunica con orden y notación correcta & Además adapta el discurso a su público con rigor & Defensa \\
  CE9 & Evita tareas o abandona ante el error & Participa de forma irregular & Colabora y aprende de sus errores & Además impulsa al equipo e incorpora la crítica & Observación, co/auto \\
\end{tablaSAlarga}

\subsection{Correspondencia con la calificación}
Los niveles se traducen a la escala numérica así: \textbf{Inicial}, 1--4; \textbf{En desarrollo}, 5--6; \textbf{Adecuado}, 7--8; \textbf{Excelente}, 9--10. En las rúbricas R1 a R4 cada nivel vale un número fijo dentro de ese tramo: 2,5; 5; 7,5 y 10. La nota de cada instrumento se calcula como indica la sección~\ref{sec:calificacion}.

\subsection{Trabajo interdisciplinar}
La tarea no se puede resolver solo con Matemáticas: requiere criterios técnicos y físicos que aportan otras materias, y cada aportación deja una evidencia evaluable.

\begin{tablaSAlarga}{colspec={Q[l,wd=86pt]X[1.4,l]X[1.4,l]X[1,l]},
  row{1}={bg=urjc, fg=white, font=\footnotesize\bfseries}, column{1}={font=\footnotesize\bfseries}}
  Materia & Qué aporta & Dónde se trabaja & Evidencia \\
  Tecnología e Ingeniería I & Energías renovables y criterios de una instalación: orientación, separación entre filas y producción & S16, con la orientación del departamento de Tecnología sobre el pliego del IDAE & Plano y cálculo de filas en la propuesta (6.1) \\
  Física y Química & Magnitudes vectoriales y modelización del rayo como dirección & S4 y S12 & Ficha F4 y modelo de sombras (7.2) \\
  Centro y comunidad & Necesidad real, destinatario de la propuesta y datos del entorno, también del centro socio & S1, S14, S16, S19 y S20 & Defensa y reflexión sobre ODS (6.2) \\
\end{tablaSAlarga}
